% Pulse area: (1) atomic law, scale, s -> boxed equation; (2) beating term by term -> closed form;
% (3) scale from the power: the three integrals; (4) C_Rabi.
% Requires: amsmath, siunitx, braket

\subsection*{Pulse area}
Start with the atomic law, then the scale of the simulation, and finally $s$ itself:
\begin{align}
\theta(\mathbf r,t_0)
&=\int_{t_0}^{t_0+T_p}\Omega(\mathbf r,t)\,dt
&&\text{definition}\\
&=C_\text{Rabi}\int_{t_0}^{t_0+T_p}I(\mathbf r,t)\,dt
&&\Omega=C_\text{Rabi}\,I\\
&=C_\text{Rabi}\,s\int_{t_0}^{t_0+T_p}\tilde I(\mathbf r,t)\,dt
&&I=s\,\tilde I\\
&=C_\text{Rabi}\,\frac{P}{\tilde P_\text{prof}}\int_{t_0}^{t_0+T_p}\tilde I(\mathbf r,t)\,dt
&&s=P/\tilde P_\text{prof}.
\end{align}
With $\tilde P_\text{prof}=\int\langle\tilde I\rangle_t\,d^2r$ this is
\begin{equation}
\boxed{\;\theta(\mathbf r,t_0)=C_\text{Rabi}\,P\;
\frac{\displaystyle\int_{t_0}^{t_0+T_p}\tilde I(\mathbf r,t)\,dt}
{\displaystyle\int\langle\tilde I\rangle_t\,d^2r}\;}
\label{eq:theta_box}
\end{equation}
This is the whole physics. Only three things enter: the atomic constant $C_\text{Rabi}$, the optical
power $P$ in the profile, and the shape of the simulated profile in the fraction. The numerator is the
light dose at the position of the atom during the pulse and contains the full beating; the
denominator is the time-averaged power of the simulated profile and only converts its arbitrary
units into watts.

\subsection*{The beating, term by term}
The time integral in the numerator can be done term by term. Inserting the Fourier series
$\tilde I=\sum_d D_d(\mathbf r)\,e^{i2\pi d f_0 t}$, each order contributes
\begin{equation}
\int_{t_0}^{t_0+T_p}e^{i2\pi d f_0 t}\,dt
=\frac{e^{i2\pi d f_0(t_0+T_p)}-e^{i2\pi d f_0 t_0}}{i\,2\pi d f_0}.
\end{equation}
Pulling out the phase of the window centre, $e^{i2\pi d f_0(t_0+T_p/2)}$, makes the bracket symmetric,
\begin{equation}
=e^{i2\pi d f_0(t_0+T_p/2)}\,\frac{e^{i\pi d f_0T_p}-e^{-i\pi d f_0T_p}}{i\,2\pi d f_0}
=e^{i2\pi d f_0(t_0+T_p/2)}\,\frac{\sin(\pi d f_0T_p)}{\pi d f_0},
\end{equation}
where $e^{ix}-e^{-ix}=2i\sin x$ was used. Writing the last factor as $T_p$ times
$\operatorname{sinc}(z)=\sin(\pi z)/(\pi z)$ gives
\begin{equation}
\int_{t_0}^{t_0+T_p}\tilde I\,dt
=T_p\sum_d D_d(\mathbf r)\,\operatorname{sinc}(d f_0T_p)\,e^{i2\pi d f_0(t_0+T_p/2)}.
\label{eq:boxcar}
\end{equation}
The three factors read directly: $T_p$ is the length of the window, the exponential says that only
the centre of the window matters for the phase, and the sinc is the Fourier transform of the window
itself, evaluated at the beat frequency $d f_0$. Integrating over a window of length $T_p$ is a
convolution with a boxcar, and in frequency space a convolution is a multiplication with the
transform of the box. The $d=0$ term is the limit $\operatorname{sinc}(0)=1$: the time average always
survives, whatever the pulse length. Because $D_{-d}=D_d^*$ the sum is real,
\begin{equation}
\int_{t_0}^{t_0+T_p}\tilde I\,dt
=T_p\Bigl[D_0(\mathbf r)+2\sum_{d>0}|D_d(\mathbf r)|\,\operatorname{sinc}(d f_0T_p)\,
\cos\!\Bigl(2\pi d f_0\bigl(t_0+\tfrac{T_p}{2}\bigr)+\arg D_d\Bigr)\Bigr].
\end{equation}
This is the same expression as for a camera exposure of duration $T_p$; a pulse and an exposure
differ only in their length. Two limits:
\begin{itemize}
\item $T_p\gg T_0$: every $d\neq0$ is damped away and $\theta\to C_\text{Rabi}\,s\,\langle\tilde I\rangle_t\,T_p$.
The pulse sees the time average.
\item $T_p\ll T_0$: $\operatorname{sinc}\to1$. The pulse does not average the beating, it samples it, and
$\theta$ depends on the trigger instant $t_0$.
\end{itemize}
\paragraph{The pulse area in closed form.}
Inserting the boxcar result \eqref{eq:boxcar} into the boxed equation \eqref{eq:theta_box} gives the
expression the code evaluates,
\begin{equation}
\theta(\mathbf r,t_0)=\frac{C_\text{Rabi}\,P\,T_p}{\tilde P_\text{prof}}
\sum_d D_d(\mathbf r)\,\operatorname{sinc}(d f_0T_p)\,e^{i2\pi d f_0(t_0+T_p/2)}.
\label{eq:theta_closed}
\end{equation}

\subsection*{The scale from the power}
The power carried by the pattern is the area integral of the time average,
\begin{equation}
P=\int\langle I(\mathbf r)\rangle_t\,d^2r=s\int\langle\tilde I(\mathbf r)\rangle_t\,d^2r\equiv s\,\tilde P_\text{prof}
\quad\Longrightarrow\quad
s=\frac{P}{\tilde P_\text{prof}}.
\end{equation}
$\tilde P_\text{prof}$ is known in closed form, and it is worth doing the three integrals explicitly.
Each spot has the field profile $g_s=A_s\,u(\mathbf r-\mathbf c_s)$ with the field amplitude
$A_s=\sqrt{a_{x,i}a_{y,j}}$ and the intensity weight $a_s\equiv A_s^2=a_{x,i}a_{y,j}$.

\paragraph{(a) The power of one spot.}
A single spot carries $\int g_s^2\,d^2r=a_sP_1$ with $P_1=\int|u|^2d^2r$. For the Gaussian,
with $x=2r^2/w_0^2$,
\begin{equation}
P_1^\text{G}=\int_0^\infty e^{-2r^2/w_0^2}\,2\pi r\,dr=\frac{\pi w_0^2}{2}\int_0^\infty e^{-x}\,dx=\frac{\pi w_0^2}{2}.
\end{equation}
For the Airy mode, with $x=kr$ and $\int_0^\infty J_\nu^2(x)\,dx/x=1/(2\nu)$,
\begin{equation}
P_1^\text{A}=\int_0^\infty\left(\frac{2J_1(kr)}{kr}\right)^2 2\pi r\,dr
=\frac{8\pi}{k^2}\int_0^\infty\frac{J_1^2(x)}{x}\,dx=\frac{4\pi}{k^2},
\qquad k=\frac{3.8317}{1.4830\,w_0}.
\end{equation}
The Airy spot carries $(4\pi/k^2)/(\pi w_0^2/2)=1.20$ times the power of a Gaussian with the same
$1/e^2$ radius, because of its rings.

\paragraph{(b) Which cross terms survive the time average.}
A pair with $k_s\neq k_{s'}$ oscillates at $(k_s-k_{s'})f_0$ and averages to zero over $T_0$, while a
degenerate pair does not oscillate at all:
\begin{equation}
\langle\tilde I(\mathbf r)\rangle_t=\sum_s g_s^2(\mathbf r)
+2\sum_{\substack{s<s'\\k_s=k_{s'}}}g_s(\mathbf r)\,g_{s'}(\mathbf r)\cos(\phi_s-\phi_{s'}).
\end{equation}
Integrating term by term over the plane gives
\begin{equation}
\tilde P_\text{prof}=P_1\sum_s a_s
+2\sum_{\substack{s<s'\\k_s=k_{s'}}}A_sA_{s'}\cos(\phi_s-\phi_{s'})\,O(d_{ss'}),
\qquad
O(d)=\int u(\mathbf r-\mathbf c_s)\,u(\mathbf r-\mathbf c_{s'})\,d^2r,
\label{eq:pprof}
\end{equation}
with $d=|\mathbf c_s-\mathbf c_{s'}|$.

\paragraph{(c) The overlap integral.}
For the Gaussian the product of the two modes is again a Gaussian, centred between the spots,
$u(\mathbf r)\,u(\mathbf r-\mathbf d)=e^{-2|\mathbf r-\mathbf d/2|^2/w_0^2}\,e^{-d^2/2w_0^2}$, so that
\begin{equation}
O^\text{G}(d)=P_1^\text{G}\,e^{-d^2/2w_0^2}.
\end{equation}
For the Airy mode the overlap is a convolution, which in Fourier space becomes a product,
$O(d)=\frac{1}{(2\pi)^2}\int|\tilde u(\mathbf q)|^2e^{i\mathbf q\cdot\mathbf d}\,d^2q$.
The Airy mode is the transform of a uniform disc, $\tilde u(\mathbf q)=\frac{4\pi}{k^2}$ for $q<k$ and
zero otherwise, so $|\tilde u|^2=\frac{4\pi}{k^2}\,\tilde u$ and
\begin{equation}
O^\text{A}(d)=\frac{4\pi}{k^2}\,\frac{1}{(2\pi)^2}\int\tilde u(\mathbf q)\,e^{i\mathbf q\cdot\mathbf d}\,d^2q
=P_1^\text{A}\,u(d).
\end{equation}
The overlap of two Airy spots is the Airy mode itself, evaluated at their distance.

\paragraph{Result for the $3\times4$ grid.}
The only degenerate pair is $(0,3)$, $(2,0)$ at $d=\SI{3.305}{\micro\metre}$ with $u(d)=0.063$; for
tone phases $0$ it adds $\SI{1.1}{\percent}$ to the power. For the phases used here
$\phi_{0,3}-\phi_{2,0}=\pi/2$, the term vanishes and
\begin{equation}
\tilde P_\text{prof}=P_1\sum_s a_s=12.70\,P_1^\text{A}=\SI{25.9}{\micro\metre\squared},
\qquad P_1^\text{A}=\SI{2.04}{\micro\metre\squared}\ (w_0=\SI{1.04}{\micro\metre}).
\end{equation}

\subsection*{The coupling constant $C_\text{Rabi}$}
A leg $i$ with intensity $I_i=\tfrac12\varepsilon_0c|E_i|^2$ couples a ground state to the excited
state $e$ with $\Omega_{i,e}=d_eE_i/\hbar$, so that
\begin{equation}
\Omega_{i,e}^2=\frac{2\,d_e^2}{\varepsilon_0c\,\hbar^2}\,I_i\equiv d_e^2\,\kappa\,I_i,
\qquad \kappa=\frac{2}{\varepsilon_0c\,\hbar^2}.
\end{equation}
The two legs share the intensity, $I_1=\beta I$ (addressing $F=2$) and $I_2=(1-\beta)I$ (addressing
$F=3$). For $|\Delta_e|\gg\Omega,\Gamma$ the excited states are eliminated adiabatically. Each of
them contributes one Raman path, and the paths add coherently:
\begin{equation}
\Omega=\sum_e\frac{\Omega_{1,e}\,\Omega_{2,e}}{2\Delta_e}
=\underbrace{\kappa\sqrt{\beta(1-\beta)}\,\left|\sum_e\frac{d_{2e}\,d_{3e}}{2\Delta_e}\right|}_{C_\text{Rabi}}\,I.
\end{equation}
The index $e$ labels a hyperfine level of the $5P$ manifold,
\begin{equation}
\ket{e}=\ket{5P_{J'},F',m'},\qquad
J'=\tfrac12\ (\mathrm{D}_1,\ F'=2,3),\qquad
J'=\tfrac32\ (\mathrm{D}_2,\ F'=1,2,3,4),
\end{equation}
so six terms for $\sigma^+$ light, where $m'=m+1$ is fixed by the polarisation.
$d_{2e}=\bra{e}\hat d_q\ket{F=2,m}$ and $d_{3e}=\bra{e}\hat d_q\ket{F=3,m}$ are the signed dipole
matrix elements, and
\begin{equation}
\Delta_e=2\pi\left[\Delta+\left(\nu_{\mathrm{D}_1}-\nu_{J'}\right)-\Delta_\text{hfs}(J',F')\right]
\end{equation}
is the detuning from level $e$ (Fig.~\ref{fig:levelscheme_detuning}): $\Delta$ is measured from the
centroid of the D$_1$ line, $\nu_{J'}$ is the centroid of the line of $e$, and $\Delta_\text{hfs}(J',F')$
the hyperfine shift of $F'$ within that line. Because both legs are two-photon resonant, both have
the same detuning from a given level $e$. The signs of the matrix elements matter, since the paths
via different $F'$ interfere.

For the clock transition $m=0$, $\beta=1/2$ and $\Delta/2\pi=+\SI{50}{\giga\hertz}$ the matrix elements
and level energies from ARC give
\begin{center}
\begin{tabular}{lrrrr}
\hline
$e$ & $d_{2e}\,(ea_0)$ & $d_{3e}\,(ea_0)$ & $\Delta_e/2\pi$ (GHz) & contribution \\ \hline
D$_1$, $F'=2$ & $-0.998$ & $0.998$ & $50.21$ & $-3.84$\\
D$_1$, $F'=3$ & $-1.411$ & $1.411$ & $49.85$ & $-7.74$\\
D$_2$, $F'=1\dots4$ & & & $\approx-7073$ & $-0.08$\\ \hline
\end{tabular}
\end{center}
with the contributions in $\si{\radian\per\second}/(\si{\watt\per\metre\squared})$, and therefore
\begin{equation}
C_\text{Rabi}=\SI{11.66}{\radian\per\second}\Big/\SI{}{\watt\per\metre\squared}.
\end{equation}
The D$_2$ line is seven terahertz away and contributes less than $\SI{1}{\percent}$.
